% ============================================================
% qp-layout.tex —— 版面模块（全品式导学件，供 main.tex \input）
% 职责：几何（A4 左/右/顶 15mm、底 20mm、twoside；headheight/headsep=0pt 撤页眉位）、宏包、
%       调色板（黑＋777777＋7A7A7A＋4D4D4D＋DDDDDD＋black!40 栏线派生档）、
%       行距/栏距/栏线/断行惩罚/浮动参数、表格默认（v4.3 总账D：深灰线色＋外框0.8pt＞内线0.4pt＋行高 pad 4.1mm）、
%       顶格化口径（\parindent=0pt）。
% 调用关系：main.tex → qp-fonts → 本模块 → qp-headfoot → qp-titles → qp-blocks。
% ============================================================
% v4.2 拍板22 冲突实害修复：bottom=20mm 正文区让位页脚带（全品 p04/p05 实测正文最低≈22.8mm、
% 块底≈11.4mm——全品正文区让开页脚带；v4.2 原 margin=15mm 实证 p1 页脚盖正文末行，故底边让 5mm）
% 字替对照 C：密度版式（任务④）——左右 margin 15mm→17.6mm（版心收窄）；bottom=20mm 让位页脚带机制不动
% 字替对照 E（0909）：margin 17.6→17.575mm（全品版心 174.85mm 实测反推，勿手抹小数）
\usepackage[a4paper,margin=17.575mm,bottom=20mm,twoside,headheight=0pt,headsep=0pt,footskip=13pt]{geometry}
% 字替对照 B/C：数学修（任务③b）——unicode-math＋TeX Gyre Termes Math（Times 风），数学变量/数字/
%   箭头换 Times 系；amssymb 与 unicode-math 冲突故卸下（body.tex 无 amssymb 专属命令，已核对）；
%   \parallel 已由 qp-fonts 的 \AtBeginDocument 在其之后再落定义（TikZ 双平行四边形复刻全品 ∥），
%   不被覆盖；原 sympx 斜体方案退役。
\usepackage{amsmath}
\usepackage{unicode-math}
\setmathfont{texgyretermes-math.otf}
% F 片C 0909 #29（数学符号间距刚性化）：unicode-math 默认 \thickmuskip=5mu plus 5mu——关系符号
%   （＝/∥ 等）两侧胶被 \emergencystretch 拉伸至自然值 2.2 倍（实测正文 1.2–2.35mm，自然
%   5mu＝0.99mm；表格单行格恒自然）。去 plus/minus 刚性化，正文墨隙回到 5mu 自然值。
\thickmuskip=5mu
\usepackage{graphicx}
\usepackage{xcolor}
\usepackage{colortbl}
\usepackage{multirow,array,calc}
\usepackage{multicol}

% ---------- 颜色（v4.3 调色板：黑＋777777＋7A7A7A＋4D4D4D＋DDDDDD＋black!40 栏线） ----------
\definecolor{midgray}{HTML}{777777}    % 灰字/括注（题侧 8pt）/章首方块深块（唯一中灰，拍板27 沿用）
\definecolor{gray122}{HTML}{7A7A7A}    % 灰122：花形右词/表线（总账A/D，全品实测 122）
\definecolor{huarule}{HTML}{4D4D4D}    % 灰77（十进制 77/255）：花形底线（总账A，全品实测 77，深灰）
\definecolor{pnumbg}{HTML}{DDDDDD}     % 页码灰块/章首浅方块（唯一浅灰）

% ---------- 版面 ----------
% 字替对照 C：栏距 7.5mm→8.8mm（目标栏宽≈83.0mm）
% 字替对照 E（0909）：栏距 8.8→9.25mm（版心 174.85mm 时栏宽＝(174.85−9.25)/2＝82.8mm，与全品栏宽实测一致；
%   C 版三表已收窄至 82.8mm，理论恰好放下——若冒 overfull 按先例最小回填留痕）
\setlength{\columnsep}{9.25mm}
\setlength{\columnseprule}{0.4pt}      % 栏间竖线（全品 p04-p06 实证宽约 0.17mm 近似 0.4pt）
% v4.3 总账I：black!25（0xBF）渲染实测灰 233 无芯偏浅 → 加深试档，取 black!40（0x99）：
% 线性折算 233×153/191≈187≈全品实测 188 带芯（出样实测见交付报告，black!35/40 双档比对取 40）
\renewcommand{\columnseprulecolor}{\color{black!40}}
\setlength{\multicolsep}{0pt}         % 章首【学习目标】标签→栏区零距（multicol 前后 skip 全局归零）
% 字替对照 C：正文 10.5pt/18pt→10.09pt/18.25pt（全品密度实测口径）
\renewcommand{\normalsize}{\fontsize{10.09pt}{18.25pt}\selectfont}
\linespread{1}                          % 行距通解：撤 ctex linestretch 1.3 叠乘，实发 18pt
\normalsize
\setlength{\parindent}{0pt}             % 顶格化（v4.1 指令5）：全品正文无首行缩进
                                        % v4.3 拍板6 例外：学习目标条目（\mubiaomu，qp-titles）悬挂 8.3mm
\clubpenalty=10000
\widowpenalty=10000
\displaywidowpenalty=10000
\tolerance=9999
% F 片C 0909 #29：\emergencystretch 8em→0pt（序言预设；multicol 进入时按 4pt×栏数重置，
%   实际生效值在 body.tex——postproc LAN_OPEN 同步归零）；\binoppenalty/\relpenalty 100→10000
%   （禁关系号后断行，消灭「a=」悬行尾；全品正文无关系号行尾实例）。
\emergencystretch=2em   % 序言预设；实际生效处在 multicols 内（multicol 进入时重置为 4pt×栏数）
\binoppenalty=10000
\relpenalty=10000
\flushbottom            % 对照§四-5（\raggedbottom 撤）
\raggedcolumns          % multicol 栏末不硬拉伸（末页参差收尾，拍板23）
\setcounter{topnumber}{4}
\setcounter{totalnumber}{4}
\setcounter{bottomnumber}{2}
\renewcommand{\topfraction}{0.95}
\renewcommand{\textfraction}{0.05}
\renewcommand{\floatpagefraction}{0.85}

% ---------- 表格默认（v4.3 总账D：内线 0.4pt（外框 0.8pt 在 kbtable 逐表组内设）；表线深灰 122
%           （\arrayrulecolor 在 kbtable 表组内设，不泄漏全局）；行高留白 pad≈4.1mm/行（\extrarowheight）；
%           tabcolsep 3pt 维持；v4.2 让位 1.05 stretch 沿用） ----------
\setlength{\arrayrulewidth}{0.4pt}
% v4.3 行高留白：全品单行行高 7.26mm（p14 逐线实测）——arraystretch 1.05 strut＋extrarowheight。
% （首版误将「pad≈4.1mm/行」直译为 extrarowheight 4.1mm，实测单行 10.45 超全品 3.2mm，0908 修正）
% v4.5-F（P-E 0909）：0.6→0.4mm＋kbtable 行尾 \\[0.45mm]——目标单行格净空 顶≈2.0/底≈2.1mm。
% （实测：m 列墨锚数学轴，extrarowheight 只长行顶 dT/dE≈+1；底净空由行尾胶补 dB/dX＝+1，探针 pe_probe.py 两轮回填）
% F 0909 收尾轮（V5 终验）：1.0→0.4mm——终验实测（_calib/pe_probe.py 360dpi 全格）E=1.0 时 m 格内容
% 系统性下坠 Δ=padT−padB≈+1.2mm（E/2 下坠），多行格不居中、单行格 顶2.75-3.14/底1.90-2.28 出窗
% （目标 2.0/2.1）；回填 0.4mm 后复测居中 ±0.3、单行格落窗（表头 8.47→7.87mm≈全品单行 7.26 档＋FakeBold 补差）
% F 片B 0909 #30（意见31 甲案标定回填）：0.4→1.7mm——甲案原定 2.6mm，实测（_tabcal 标定矩阵 0909，
% 表1/表2 真内容 600dpi ink）padT 对 E 的响应 dT/dE≈+1.09：E=2.6 时多行格顶净空 4.5mm 出窗
% （目标 3.5±0.3）；E=1.7 时 padT 3.4–3.5 落窗、表头行高 9.18（全品 9.03）落 9.0±0.3 窗。底侧净空由
% 行尾胶经 \tabrowglue 按行折数分档供给（见下）——E=1.7 为顶/表头双窗共解值（E≥1.8 表头出窗）。
\setlength{\extrarowheight}{1.7mm}
% 表行尾胶自动分档（F 片B 0909 #30）：\\[\tabrowglue]——m 列行尾胶在折数 ≥2 时被居中机制吸收，
% 底净空需超过阈值才生效（实测阈值≈2.44mm、超过后 dB/dG≈1）；折数用 TeX 自测（\prevgraf）：
% 每数据格首格 \rowreset 清零、格末 \marklines 记 max 行数，行尾胶按 1/2/≥3 折取 0.85/4.65/7.60mm
% ——实测多行格行级最紧侧 顶3.1–3.5/底3.0–4.0（基线 0.9–2.2），全品 p04 多行格 3.46–4.59 档内；
% 单行格（表2 共线行）底净空保 2.2（原 2.1/2.2 档），行高 7.82→9.58 随 E 全局垫顶（登记）。
\newcount\rowmax
\newcount\celllines
\newcommand{\rowreset}{\global\rowmax=0 }
\newcommand{\marklines}{\par\celllines=\prevgraf\relax
  \ifnum\celllines>\rowmax \global\rowmax=\celllines \fi}
\newcommand{\tabrowglue}{\ifnum\rowmax>2 7.60mm\else\ifnum\rowmax>1 4.65mm\else 0.85mm\fi\fi}
\setlength{\tabcolsep}{3pt}
\renewcommand{\arraystretch}{1.05}
\renewcommand{\multirowsetup}{\centering}
